\section{A finite slice of the support}
\label{sec:support}

We complete the proof by applying the slicing theorem to the support of
$F$.  First we arrange that a nonzero generic stalk specializes to the
closed fiber.  We then cut its support to a finite scheme over the trait.
On such a scheme, geometric nearby cycles are a finite direct sum of stalk
complexes, so they cannot vanish if one of those stalks is nonzero.

\subsection{Bringing a generic stalk to the closed fiber}

We use the bounded Hecke spaces of Section~\ref{sec:hecke} to extend
bundles.  The argument applies to connected reductive groups, including
those with a central torus.

\begin{lemma}
\label{support:extension}
Let $\mathscr X\to\Spec R$ be the smooth proper curve fixed in
Section~\ref{sec:specialization}, and let $K$ be the fixed algebraic
closure of $\operatorname{Frac}(R)$.  Every $G$-bundle on $\mathscr X_K$
extends to a $G$-bundle on $\mathscr X_{R'}$ for some finite extension
$\operatorname{Frac}(R')/\operatorname{Frac}(R)$ inside $K$, where $R'$
is the integral closure of $R$ in that extension.
\end{lemma}

\begin{proof}
Let $\mathcal P_K$ be the bundle.  By
\cite[Theorem~2 and Remark~2(b)]{DrinfeldSimpson}, it is trivial on a
nonempty open subset of $\mathscr X_K$.  Here Theorem~2 applies to
connected reductive groups, and the base field $K$ is algebraically
closed.  Choose a trivialization away from distinct points
$x_1,\ldots,x_r\in\mathscr X(K)$.  Gluing at these points gives a chain
of one-point modifications
\[
 \mathcal P_0\dashrightarrow\mathcal P_1\dashrightarrow\cdots
 \dashrightarrow\mathcal P_r=\mathcal P_K,
 \qquad \mathcal P_0\text{ trivial},
\]
whose $i$th modification has leg $x_i$.  Choose a Schubert bound for each
of the finitely many relative positions in this chain.

Starting with the trivial bundle over $S=\Spec R$, form the space of
chains satisfying these bounds, allowing the legs to vary on
$\mathscr X$.  At each stage the bounded Hecke space is representable
and proper over the preceding bundle and its new leg.  Since
$\mathscr X\to S$ is proper, induction shows that the space of chains is
a proper algebraic space of finite type over $S$.  It carries the
universal final bundle and contains the chosen chain as a $K$-point.
This point descends to a finite extension of $\operatorname{Frac}(R)$:
the space is of finite type and $K$ is algebraic over that field.  The
valuative criterion extends the descended point over $R'$, and the
universal final bundle gives the required extension.  The construction
allows the legs to coincide on the closed fiber, since it uses successive
one-point modifications.
\end{proof}

Suppose now that a spectral component is missing, and let $F$ be the
nonzero compact object supplied by Proposition~\ref{spec:missing}.
It is constructible on finite-type smooth charts.  There is therefore a
$K$-valued bundle at which its stalk is nonzero: a nonzero constructible
complex on a finite-type scheme over $K$ has a nonzero stalk at a closed
$K$-point.  Apply Lemma~\ref{support:extension} to this bundle and replace
$R$ by the resulting $R'$.  All specialization statements from
Section~\ref{sec:specialization} remain valid with the same geometric
generic field $K$.

Choose a finite-type smooth affine chart $b:U\to\mathcal B$ through
the special value of the extended bundle.  The fiber product of this
chart with its trait section is smooth over $S$ and has a $k$-point.
It has an $R$-point lifting that $k$-point.  Indeed, after taking an
affine \'etale neighborhood if needed, formal smoothness gives
compatible lifts modulo every power of the uniformizer, and completeness
of $R$ gives an $R$-point of the affine scheme.  Thus the trait section
lifts to $U$.  After shrinking around its closed point, we may suppose
that $U/S$ has constant relative dimension $N$.  The generic stalk of
$F_U=b_K^*F$ on the lifted section is nonzero, so its support has closure
meeting $U_s$.

\subsection{A slice with nonzero nearby cycles}

The remaining argument uses only the dimension of a cone in the
special-fiber cotangent bundle.  We state it separately to make clear
that the support strata, rather than the sheaf, are the objects descended
to a finite field extension.

\begin{lemma}
\label{support:finite-slice}
Let $R$ be a complete discrete valuation ring with algebraically closed
residue field $k$ of characteristic different from $\ell$, let
$S=\Spec R$, and let $K$ be an algebraic closure of its fraction field.  Let $U/S$ be a smooth affine
scheme of constant
relative dimension $N$, let $\mathcal F$ be a bounded constructible
$E$-complex on $U_K$, and let
$\Lambda\subset T^*U_s$ be a closed conical subset with
$\dim\Lambda\leq N$.
Suppose that the closure in $U$ of the nonzero-stalk locus of
$\mathcal F$ meets $U_s$.
Then, after a finite extension of the trait and shrinking $U$, there
exist $u\in U_s(k)$, an integer $0\leq d\leq N$, and functions
$f_1,\ldots,f_d$ vanishing at $u$ such that their special-fiber
differentials are independent,
\[
 \operatorname{span}\{d(f_{1,s})_u,\ldots,d(f_{d,s})_u\}
       \cap\Lambda_u\subseteq\{0\},
\]
and, for $J=V(f_1,\ldots,f_d)$,
\begin{equation}
\label{eq:support-nonzero}
 \bigl(\Psi_J(\mathcal F|_{J_K})\bigr)_u\ne0.
\end{equation}
For $d=0$, the list of functions is empty and $J=U$.
\end{lemma}

\begin{proof}
\emph{Descend the finite support data.}
Let $A_1,\ldots,A_t$ be the irreducible components of the closure of
the nonzero-stalk locus in $U_K$.  Bounded constructibility gives a
proper closed subset $B_i\subsetneq A_i$ such that every stalk on
$A_i\setminus B_i$ is nonzero.  Give these closed sets their reduced
structures.  Their finitely generated ideals descend to a common finite
extension of $\operatorname{Frac}(R)$; enlarge that extension to descend
all containments as well.  Replace the trait accordingly.  The new
support closure still meets
the special fiber: the finite map from the new model to the old one
maps its closure onto the old closure, since it is closed and agrees
on the geometric generic support.  This step does not assert that
$\mathcal F$ descends.

Write $\overline A_i$ and $\overline B_i$ for their horizontal closures
in $U$, and put $e_i=\dim A_i$.  Each $\overline A_i$ is integral and
dominates $S$.  Every nonempty special fiber
$(\overline A_i)_s$ is pure of dimension $e_i$
\cite[Tag~0B2J]{Stacks}.  Applying the same statement to the
irreducible components of $B_i$ shows that
\[
 \dim(\overline B_i)_s<e_i
 \quad\text{whenever }(\overline B_i)_s\ne\varnothing.
\]
The trait is excellent, and hence universally catenary: a complete
Noetherian local ring is excellent \cite[Tag~07QW]{Stacks}.
Consequently, at every closed point $v\in(\overline A_i)_s$, the
dimension formula gives
\begin{equation}
\label{eq:support-local-dimension}
 \dim\mathcal O_{\overline A_i,v}=e_i+1
 \qquad\text{\cite[Tag~02JT]{Stacks}}.
\end{equation}

\emph{Choose the special point.}
Let $d$ be maximal among the $e_i$ for which
$(\overline A_i)_s$ is nonempty.  Such an index exists by the support
assumption.  Choose a reduced irreducible component $T$ of dimension
$d$ in one of these special fibers.  At a general closed point of $T$,
its reduced structure is smooth, no other distinct reduced component of
$\bigcup_i(\overline A_i)_s$ passes through the point, and none of the
$\overline B_i$ passes through it.  To see the last assertion, a
component with empty special fiber contributes nothing, while every
remaining bad special locus has dimension strictly below $d$.
Different horizontal components $\overline A_i$ may have the same
reduced special component $T$; we retain all of them.

We can also arrange that
\begin{equation}
\label{eq:support-cone-fiber}
 \dim\Lambda_u\leq N-d.
\end{equation}
Indeed, restrict $\Lambda$ to $T$.  Its components that do not dominate
$T$ can be excluded over a proper closed subset of $T$.  The generic
fiber of each remaining component has dimension at most $N-d$, since
$\dim\Lambda\leq N$.  The same bound holds over a nonempty open subset
of $T$.  Choose $u\in T(k)$ satisfying all these conditions.

\emph{Choose transverse equations.}
Let $V=T_u^*U_s$, so $\dim V=N$.  If $d>0$, a general $d$-plane
$L\subset V$ satisfies $L\cap\Lambda_u\subseteq\{0\}$.  Indeed,
\eqref{eq:support-cone-fiber} gives
$\dim\mathbb P(\Lambda_u)\leq N-d-1$, and the incidence of
$d$-planes meeting this projective set is a proper closed subset of
$\Gr(d,V)$.  The condition that
\[
 L\longrightarrow T_u^*T
\]
be an isomorphism is another nonempty open condition.  These two open
subsets meet because the Grassmannian is irreducible.  Choose such an
$L$ and a basis of it.  For $d=0$, take $L=0$.

Lift the basis to functions $f_1,\ldots,f_d$ vanishing at $u$, on an
affine neighborhood in $U$.  Their restrictions to the smooth
$d$-dimensional variety $T$ form a regular system of parameters at $u$.
Hence, for $J=V(f_1,\ldots,f_d)$ and
\[
 Z=J\cap\bigcup_i\overline A_i,
\]
the point $u$ is isolated in $Z_s$ after shrinking.  This assertion
concerns the union of all the support closures: its reduced special
fiber near $u$ is precisely $T$, even if several horizontal components
specialize to $T$.

We have obtained the required differentials.  It remains to show that
the slice retains a nonzero generic stalk and that this stalk survives
specialization.

\emph{Find a generic generization in the slice.}
Choose an index $i$ for which $e_i=d$ and $T$ is a component of
$(\overline A_i)_s$.  By \eqref{eq:support-local-dimension}, the local
ring of $\overline A_i$ at $u$ has dimension $d+1$.  Quotienting by
$f_1,\ldots,f_d$ leaves dimension at least one.  Its quotient by a
uniformizer $\pi$ has dimension zero, by the isolation just proved.
Thus some prime of the cut local ring does not contain $\pi$:
otherwise $\pi$ would be nilpotent and the ring would have dimension
zero.  This gives a generic-fiber point $z$ specializing to $u$.

The morphism $Z\to S$ is quasi-finite at $u$, since $u$ is an isolated
closed point of its fiber \cite[Tag~01TH]{Stacks}.  Its quasi-finite
locus is open, so it contains $z$.  Thus $z$ has residue field finite
over $\operatorname{Frac}(R)$.  Its geometric lifts to $K$ avoid
$B_i$: a point in the closed set $\overline B_i$ could not specialize
to $u\notin\overline B_i$.  Each such lift therefore has a nonzero
stalk of $\mathcal F|_{J_K}$.

\emph{Compute nearby cycles on a finite neighborhood.}
Take an affine neighborhood of $u$ inside the quasi-finite locus of
$Z\to S$.  The henselian decomposition of a finite-type algebra
\cite[Tag~04GJ]{Stacks} provides a finite local factor containing $u$.
It defines an open neighborhood $Q$ of $u$ in $Z$, finite over $S$,
whose special fiber has only the point $u$.  Equivalently, the
separated form of Zariski's Main Theorem embeds the quasi-finite
neighborhood into a finite $S$-scheme; its henselian local factor at
$u$ lies in the original open because an open subset of a local
spectrum containing the closed point is the entire spectrum.
Every generization of $u$, including $z$, lies in $Q$.

The restriction $\mathcal F|_{J_K}$ is supported on the closed subset
$Z_K$.  To use the finite neighborhood, remove $Z\setminus Q$ from
$J$.  We realize this removal in the ambient chart: the complement is
closed in $J$, hence in $U$, so remove it from $U$ and then choose an
affine neighborhood of $u$.  This neighborhood contains all of $Q$,
since every point of the finite local scheme $Q$ specializes to $u$.
The new $J$ is still cut out by the same functions in a smooth affine
chart; the change leaves its nearby-cycle stalk at $u$ unchanged and
makes $Q$ closed in $J$.  If $\iota:Q\hookrightarrow J$ denotes this
closed immersion and $\mathcal H=\mathcal F|_{Q_K}$, localization gives
\[
 \mathcal F|_{J_K}\simeq\iota_{K,*}\mathcal H.
\]
Proper compatibility for $\iota$ therefore reduces
\eqref{eq:support-nonzero} to the stalk of $\Psi_Q\mathcal H$ at $u$.
Let $a:Q\to S$ be the finite structural map.  Its special fiber has
only the point $u$, so proper compatibility gives
\[
 (\Psi_Q\mathcal H)_u
 \simeq \Psi_S\bigl(\RGamma(Q_K,\mathcal H)\bigr)
 \simeq \RGamma(Q_K,\mathcal H).
\]
The second identification uses geometric specialization on the trait,
which is the identity on coefficient complexes.  The finite $K$-scheme
$Q_K$ has the \'etale topos of a finite discrete set.  Consequently
\[
 \RGamma(Q_K,\mathcal H)
   \simeq\bigoplus_{x\in|Q_K|}\mathcal H_x.
\]
At least one summand is nonzero, as proved above, so this complex is
nonzero.  This proves \eqref{eq:support-nonzero}.

These are geometric nearby cycles with the monodromy action forgotten;
no inertia invariants are taken.  The displayed finite direct-sum
calculation is compatible with finite coefficients, their
$\ell$-adic limit, inversion of $\ell$, and extension to $E$.
It therefore applies to the specialization functor of
Section~\ref{sec:specialization} without requiring a field of descent
for $\mathcal F$.
\end{proof}

\subsection{Completion of the full-support proof}

\begin{proof}[Proof of Theorem~\ref{thm:full-support}]
If $Y'\ne Y$, Proposition~\ref{spec:missing} supplies the nonzero
compact object $F$ used above.  Lemma~\ref{support:extension} and the
subsequent chart construction give a smooth affine $U\to\mathcal B$
of constant relative dimension $N$ for which the closure of the
nonzero-stalk locus of $F_U$ meets $U_s$.

Let $\mathcal N_U\subset T^*U_s$ be the pullback of the global
nilpotent cone.  Under either displayed hypothesis regime, the
nilpotent-parabolic assumption and invariant form give the
half-dimensionality theorem of
\cite[Appendix~D.1.1--D.1.8]{AGKRRV}.  In its smooth-chart formulation,
\cite[Appendix~D.1.2]{AGKRRV}, this says exactly that
\[
 \dim\mathcal N_U\leq\dim U_s=N.
\]
Apply Lemma~\ref{support:finite-slice} with
$\mathcal F=F_U$ and $\Lambda=\mathcal N_U$.  It gives a slice
$J=V(f_1,\ldots,f_d)$ and $u\in J_s(k)$ for which
\[
 \bigl(\Psi_J(F_U|_{J_K})\bigr)_u\ne0,
\]
while the independent differentials of its equations span a plane
meeting $(\mathcal N_U)_u$ only at zero.  The higher-codimension
vanishing of Proposition~\ref{slice:higher} applies to this same chart
and slice and says that this stalk is zero.  This contradiction also
covers $d=0$, when the test is the original smooth-chart vanishing.

Thus the complementary union is empty.  The primed equivalence recalled
in Section~\ref{sec:specialization} is consequently the full restricted
equivalence asserted in Theorem~\ref{thm:full-support}.
\end{proof}
